\documentclass[aspectratio=169, 11pt]{beamer}

\usepackage[utf8]{inputenc}
\usepackage[T1]{fontenc}
\usepackage[brazil]{babel}
\usepackage{graphicx}
\usepackage{booktabs}
\usepackage{amsmath}
\usepackage{tikz}

\usetheme{Madrid}
\usecolortheme{default}

\definecolor{Canopy}{HTML}{1F5C3A}
\definecolor{Leaf}{HTML}{4F8A3D}
\definecolor{Moss}{HTML}{6B7D35}
\definecolor{Fern}{HTML}{DDEBD4}
\definecolor{Mist}{HTML}{F6FAF1}
\definecolor{Bark}{HTML}{3A2E24}
\definecolor{Earth}{HTML}{8A6A3D}

\setbeamercolor{normal text}{fg=Bark,bg=Mist}
\setbeamercolor{structure}{fg=Canopy}
\setbeamercolor{title}{fg=Canopy}
\setbeamercolor{subtitle}{fg=Moss}
\setbeamercolor{author}{fg=Bark}
\setbeamercolor{institute}{fg=Bark}
\setbeamercolor{date}{fg=Bark}
\setbeamercolor{frametitle}{fg=white,bg=Canopy}
\setbeamercolor{title in head/foot}{fg=white,bg=Canopy}
\setbeamercolor{author in head/foot}{fg=white,bg=Moss}
\setbeamercolor{date in head/foot}{fg=white,bg=Canopy}
\setbeamercolor{title page title}{fg=white,bg=Canopy}
\setbeamercolor{title page subtitle}{fg=Canopy,bg=Mist}
\setbeamercolor{block title}{fg=white,bg=Canopy}
\setbeamercolor{block body}{fg=Bark,bg=Fern}
\setbeamercolor{item}{fg=Leaf}
\setbeamercolor{subitem}{fg=Moss}

\setbeamertemplate{navigation symbols}{}
\setbeamertemplate{itemize item}{\color{Leaf}$\blacktriangleright$}
\setbeamertemplate{itemize subitem}{\color{Moss}$\bullet$}
\setbeamertemplate{enumerate item}{\color{Canopy}\insertenumlabel.}
\setbeamertemplate{frametitle}{
    \vspace{0.15cm}
    \begin{beamercolorbox}[wd=\paperwidth,ht=0.72cm,dp=0.22cm,leftskip=0.55cm]{frametitle}
        \usebeamerfont{frametitle}\insertframetitle
    \end{beamercolorbox}
}
\setbeamertemplate{footline}{
    \leavevmode%
    \hbox{%
    \begin{beamercolorbox}[wd=.72\paperwidth,ht=2.4ex,dp=1ex,leftskip=.35cm]{author in head/foot}%
        Grupo 10 · \textit{Cecropia pachystachya}
    \end{beamercolorbox}%
    \begin{beamercolorbox}[wd=.28\paperwidth,ht=2.4ex,dp=1ex,center]{date in head/foot}%
        \insertframenumber/\inserttotalframenumber
    \end{beamercolorbox}}%
}

% Slide divisor de parte.
\newcommand{\partdivider}[2]{%
    \begin{frame}[plain]
        \begin{tikzpicture}[remember picture,overlay]
            \fill[Canopy] (current page.south west) rectangle (current page.north east);
            \fill[Leaf] ([xshift=12.2cm,yshift=-1.6cm]current page.south west) circle (2.6cm);
            \fill[Fern] ([xshift=12.9cm,yshift=-1.1cm]current page.south west) circle (1.7cm);
            \fill[Moss] ([xshift=-0.7cm,yshift=1.3cm]current page.south west) circle (2.1cm);
        \end{tikzpicture}
        \vfill
        \begin{center}
            {\normalsize\color{Fern} #1}\\[0.5em]
            {\LARGE\bfseries\color{white} #2}
        \end{center}
        \vfill
    \end{frame}
}

\title[Detecção de Embaúba]{Detecção e Segmentação de Copas de \textit{Cecropia} (Embaúba) em Ortomosaico de VANT}
\subtitle{Duas abordagens: Visão Computacional Clássica e Redes Neurais (SAM 3 + YOLOv8)}
\author[Grupo 10]{\texorpdfstring{Grupo 10\\
\small Augusto de H. V. Guerner · Eduardo G. de Lazari · Fábio H. A. Coelho\\
\small Gabriel Arantes F. Gualda · Nena A. Impanta}{Grupo 10}}
\institute{INE410159 — Visão Computacional \\ UFSC}
\date{Julho de 2026}
\hypersetup{
    pdftitle={Detecção e Segmentação de Copas de Cecropia em Ortomosaico de VANT},
    pdfauthor={Grupo 10}
}

\begin{document}

\begin{frame}[plain]
    \begin{tikzpicture}[remember picture,overlay]
        \fill[Mist] (current page.south west) rectangle (current page.north east);
        \fill[Canopy] (current page.north west) rectangle ([yshift=-2.35cm]current page.north east);
        \fill[Leaf] ([xshift=10.8cm,yshift=-7.5cm]current page.north west) circle (3.2cm);
        \fill[Fern] ([xshift=11.4cm,yshift=-7.0cm]current page.north west) circle (2.3cm);
        \fill[Moss] ([xshift=-0.8cm,yshift=0.4cm]current page.south west) circle (2.6cm);
    \end{tikzpicture}
    \vspace{0.35cm}
    \begin{beamercolorbox}[wd=\textwidth,center,sep=0.3cm]{title page title}
        {\LARGE\bfseries Detecção e Segmentação de Copas de \textit{Cecropia} (Embaúba) em Ortomosaico de VANT}
    \end{beamercolorbox}

    \vspace{0.45cm}
    \begin{center}
        {\large\color{Canopy}\bfseries Visão Computacional Clássica \& Redes Neurais (SAM 3 + YOLOv8)}

        \vspace{0.55cm}
        {\Large\color{Bark}\bfseries Grupo 10}

        \vspace{0.25cm}
        {\small Augusto de H. V. Guerner · Eduardo G. de Lazari · Fábio H. A. Coelho}\\
        {\small Gabriel Arantes F. Gualda · Nena A. Impanta}

        \vspace{0.55cm}
        {\small INE410159 — Visão Computacional}\\
        {\small UFSC · Julho de 2026}
    \end{center}
\end{frame}

\begin{frame}{Contexto do Trabalho}
    \begin{block}{Duas frentes complementares sobre o \textbf{mesmo} ortomosaico}
        \begin{itemize}
            \item \textbf{Visão clássica}: HSV, morfologia, contornos e filtros geométricos.
            \item \textbf{SAM 3 e YOLOv8}: detecção e segmentação com modelos modernos.
        \end{itemize}
    \end{block}

    \vspace{0.4em}
    Esta apresentação cobre as \textbf{duas frentes} para o mapeamento de copas de
    \textit{Cecropia pachystachya} em ortomosaico de VANT do experimento MataDIV,
    e as compara ao final.
\end{frame}

\begin{frame}{Problema Biológico e Objetivo}
    \begin{itemize}
        \item Detectar e segmentar copas de \textit{Cecropia pachystachya} em mosaico aéreo de alta resolução.
        \item Espécie pioneira, indicadora de regeneração, com copa umbeliforme distinta.
        \item Comparar uma abordagem \textbf{clássica} (regras de cor e forma) com uma \textbf{neural} (padrões aprendidos).
        \item Reduzir o esforço manual e os falsos positivos causados por vegetação de cor semelhante.
    \end{itemize}
\end{frame}

\begin{frame}{Área de Estudo e Ortomosaico (comum às duas frentes)}
    \begin{columns}
        \begin{column}{0.52\textwidth}
            \begin{itemize}
                \item Experimento \textbf{MataDIV} — Itatinga, SP (ESALQ/USP).
                \item 144 parcelas, $\sim$8 ha, 6 espécies nativas + alta diversidade.
                \item Sobrevoo: DJI Matrice 4T, 13/03/2026, 23 m de altitude.
                \item GSD $\approx$ 0,85 cm/pixel.
            \end{itemize}
        \end{column}
        \begin{column}{0.48\textwidth}
            \begin{block}{Ortomosaico}
            \small
            48.264 $\times$ 47.029 px · 4 bandas (RGBA) · $\sim$8,46 GB\\[0.3em]
            SIRGAS 2000 / UTM 22S (EPSG:31982)
            \end{block}
        \end{column}
    \end{columns}
\end{frame}

% ═════════════════════════════════════════════════════════════════════════════
\partdivider{Parte I}{Solução com Visão Computacional Clássica}
% ═════════════════════════════════════════════════════════════════════════════

\begin{frame}{Recorte e Amostras (Abordagem Clássica)}
    \begin{table}
        \centering
        \begin{tabular}{ll}
            \toprule
            \textbf{Parâmetro} & \textbf{Valor} \\
            \midrule
            Tiles & 992 JPEGs de 2048 $\times$ 2048 px \\
            CRS & EPSG:31982 (UTM Zona 22S) \\
            Resolução espacial & $\approx$ 0,00848 u/px \\
            Tamanho em disco & 1,4 GB \\
            Tiles pretos ignorados & 402 \\
            Tiles processados & 590 \\
            Validação manual & 50 tiles revisados \\
            \bottomrule
        \end{tabular}
    \end{table}
\end{frame}

\begin{frame}{Motivação da Pipeline Clássica}
    \begin{itemize}
        \item A embaúba apresenta copa com tonalidade verde-clara/amarelada no mosaico.
        \item A segmentação HSV gera candidatos por cor.
        \item A morfologia conecta fragmentos e remove ruído.
        \item Os contornos transformam a máscara em objetos mensuráveis.
        \item Filtros geométricos reduzem candidatos incompatíveis com copas plausíveis.
    \end{itemize}
\end{frame}

\begin{frame}{Visão Geral do Pipeline}
    \centering
    \includegraphics[height=0.82\textheight]{../../data/output/tiles/tile_0681/grid.png}
\end{frame}

\begin{frame}{Etapas do Pipeline}
    \begin{enumerate}
        \item Gaussian Blur: kernel $9 \times 9$
        \item Conversão BGR $\rightarrow$ HSV
        \item Limiarização HSV: H[44,58] S[144,231] V[104,203]
        \item Fechamento morfológico: elipse $25 \times 25$
        \item Abertura morfológica: elipse $9 \times 9$
        \item Detecção de contornos externos
        \item Filtro de área mínima: área $> 10.000$ px²
        \item Filtro final: área/circularidade + área máxima + solidez
        \item Convex Hull para visualização e exportação
    \end{enumerate}
\end{frame}

\begin{frame}{Filtro Final de Copa}
    \begin{block}{Regra de decisão}
    Um candidato é aceito quando:

    \[
        (\text{área} \geq 33.798 \;\; \text{ou} \;\; \text{circularidade} \leq 0{,}1551)
    \]

    e também:

    \[
        \text{área} \leq 600.000
        \qquad
        \text{solidez} \geq 0{,}48
    \]
    \end{block}

    \begin{itemize}
        \item Área/circularidade separa copas grandes ou irregulares.
        \item Área máxima remove blobs fundidos gigantes.
        \item Solidez remove regiões excessivamente fragmentadas.
    \end{itemize}
\end{frame}

\begin{frame}{Resultados Gerais}
    \begin{table}
        \centering
        \begin{tabular}{ll}
            \toprule
            \textbf{Métrica} & \textbf{Valor} \\
            \midrule
            Tiles processados & 590 \\
            Tiles com detecção & 507 (85,9\%) \\
            Total de detecções & 2.003 \\
            Detecções por tile & $3{,}39 \pm 2{,}69$ \\
            Área total estimada & 13.164,42 m² \\
            Área média por região & 91.406 px² ($\approx$ 6,57 m²) \\
            Área máxima por região & 599.406 px² ($\approx$ 43,10 m²) \\
            Cobertura HSV média & 13,862\% \\
            Tempo total & 215,6 s \\
            Tempo médio por tile & $400{,}8 \pm 163{,}6$ ms \\
            \bottomrule
        \end{tabular}
    \end{table}
\end{frame}

\begin{frame}{Distribuições Espaciais}
    \begin{columns}
        \begin{column}{0.5\textwidth}
            \centering
            \includegraphics[width=\textwidth]{../../data/output/hist_deteccoes_por_tile.png}
            \small Detecções por tile
        \end{column}
        \begin{column}{0.5\textwidth}
            \centering
            \includegraphics[width=\textwidth]{../../data/output/hist_areas_deteccao.png}
            \small Áreas das detecções
        \end{column}
    \end{columns}
\end{frame}

\begin{frame}{Cobertura de Vegetação e Top 10}
    \begin{columns}
        \begin{column}{0.5\textwidth}
            \centering
            \includegraphics[width=\textwidth]{../../data/output/cobertura_percentual.png}
            \small Cobertura HSV por tile
        \end{column}
        \begin{column}{0.5\textwidth}
            \centering
            \includegraphics[width=\textwidth]{../../data/output/top10_tiles.png}
            \small Top 10 por detecções
        \end{column}
    \end{columns}
\end{frame}

\begin{frame}{Validação Manual}
    \begin{table}
        \centering
        \begin{tabular}{lr}
            \toprule
            \textbf{Métrica} & \textbf{Valor} \\
            \midrule
            Tiles revisados & 50/50 \\
            Detecções avaliadas & 207 \\
            Verdadeiros positivos (TP) & 179 \\
            Falsos positivos (FP) & 28 \\
            Falsos negativos / faltantes (FN) & 79 \\
            Precisão & 86,5\% \\
            Recall & 69,4\% \\
            F1 & 77,0\% \\
            \bottomrule
        \end{tabular}
    \end{table}

    \vspace{0.5em}
    \small Alta precisão indica poucos falsos positivos na amostra revisada; o recall mostra que ainda há copas reais fora da saída final.
\end{frame}

\begin{frame}{Exemplos da Validação}
    \centering
    \includegraphics[height=0.78\textheight]{../../data/output/exemplos_validacao_tp_fp_fn.png}

    \vspace{0.3em}
    \small Saída final do detector e exemplos de TP, FP e FN.
\end{frame}

\begin{frame}{Desempenho Computacional}
    \begin{columns}
        \begin{column}{0.5\textwidth}
            \centering
            \includegraphics[width=\textwidth]{../../data/output/tempo_processamento.png}
            \small Tempo por tile
        \end{column}
        \begin{column}{0.5\textwidth}
            \centering
            \includegraphics[width=\textwidth]{../../data/output/memoria_por_tile.png}
            \small Memória por tile
        \end{column}
    \end{columns}

    \vspace{0.8em}
    \small Tempo médio de 400,8 ms/tile e memória próxima de 64 MB/tile.
\end{frame}

\begin{frame}{Limitações da Abordagem Clássica}
    \begin{itemize}
        \item Confusão espectral: outras espécies podem cair na mesma faixa HSV.
        \item Sensibilidade a iluminação e sombra.
        \item Possível fusão de copas próximas nas etapas morfológicas.
        \item Possível duplicação por sobreposição entre tiles.
        \item Recall limitado: parte das copas reais não vira candidato HSV ou é filtrada.
    \end{itemize}
\end{frame}

% ═════════════════════════════════════════════════════════════════════════════
\partdivider{Parte II}{Solução com Redes Neurais (SAM 3 + YOLOv8)}
% ═════════════════════════════════════════════════════════════════════════════

\begin{frame}{Motivação: SAM 3 e Curadoria por Exclusão}
    \begin{itemize}
        \item \textbf{SAM 3} (Meta AI): modelo fundacional treinado em $>$1 bilhão de máscaras.
        \item Com \textit{prompt} textual ``tree'', segmenta apenas copas de árvores.
        \item \textbf{Curadoria por exclusão}: em vez de anotar do zero, o operador
        \textit{remove} os segmentos que não são a espécie-alvo.
        \item Reduz drasticamente o esforço de anotação e a barreira de entrada do
        \textit{deep learning}.
    \end{itemize}
\end{frame}

\begin{frame}{Pipeline Semi-automatizado em 6 Etapas}
    \begin{enumerate}
        \item Geração do dataset com SAM 3 (\textit{text prompt} ``tree'') + curadoria.
        \item Divisão espacial do dataset (evita \textit{data leakage}).
        \item Correção semi-automatizada das anotações.
        \item Treinamento do YOLOv8x (otimização bayesiana).
        \item Inferência e pós-processamento (duas rodadas + deduplicação).
        \item Segmentação das copas novas com SAM 3 (\textit{prompt} geométrico).
    \end{enumerate}
\end{frame}

\begin{frame}{Etapa 1 — Geração do Dataset}
    \begin{itemize}
        \item Ortomosaico recortado em tiles de 4.096 px (overlap 1.024) $\to$ SAM 3.
        \item 4.569 polígonos gerados $\to$ curadoria em QGIS $\to$ \textbf{1.109 máscaras} de \textit{Cecropia}.
        \item \textit{Bounding boxes} geradas automaticamente das máscaras (sem traçado manual).
        \item Tiles de treino de 896 px (overlap 50\%): 1.109 indivíduos $\to$ \textbf{7.488 instâncias}.
    \end{itemize}
    \vspace{0.4em}
    \small As máscaras já servem de segmentação final para os 975 indivíduos dentro do MataDIV.
\end{frame}

\begin{frame}{Etapas 2 e 3 — Divisão Espacial e Correção}
    \begin{itemize}
        \item \textbf{Divisão espacial}: clusterização de tiles por grafo (BFS) + Monte Carlo (70/15/15) para evitar vazamento de dados.
        \item \textbf{Correção semi-automatizada}: modelo aponta detecções sem correspondência (IoU $<$ 0,5), operador confirma copas reais não anotadas.
    \end{itemize}
    \begin{table}
        \centering
        \small
        \begin{tabular}{lrrr}
            \toprule
            \textbf{Conjunto} & \textbf{Imagens} & \textbf{Inst. (orig.)} & \textbf{Inst. (corr.)} \\
            \midrule
            Treinamento & 776   & 6.300 & 6.300 \\
            Validação   & 167   & 543   & 687 (+144) \\
            Teste       & 166   & 645   & 799 (+154) \\
            \midrule
            Total       & 1.109 & 7.488 & 7.786 \\
            \bottomrule
        \end{tabular}
    \end{table}
\end{frame}

\begin{frame}{Etapa 4 — Treinamento do YOLOv8x}
    \begin{itemize}
        \item Variante de maior capacidade (Ultralytics); só detecção — segmentação fica com o SAM 3.
        \item \textbf{Otimização bayesiana} (Optuna / TPE): 21 hiperparâmetros, 30 tentativas $\times$ 20 épocas.
        \item Três categorias: otimizador, pesos da perda e \textit{data augmentation}.
        \item Importância dos hiperparâmetros via \textit{Random Forest} (Mean Decrease Impurity).
        \item Treino final: 500 épocas, imagem 896 px, \textit{batch} 8, AdamW.
    \end{itemize}
\end{frame}

\begin{frame}{Etapas 5 e 6 — Inferência e Segmentação}
    \begin{columns}
        \begin{column}{0.5\textwidth}
            \textbf{Inferência (pós-processamento)}
            \begin{itemize}
                \small
                \item Duas rodadas (grade + offset 50\%).
                \item Descarte de bordas e deduplicação.
                \item Filtros: confiança $\geq$ 0,65; área 1,75–40 m².
                \item Cruzamento com máscaras SAM 3 (3 categorias).
            \end{itemize}
        \end{column}
        \begin{column}{0.5\textwidth}
            \textbf{Segmentação SAM 3 (\textit{bbox prompt})}
            \begin{itemize}
                \small
                \item Recorte por bbox + \textit{padding} 300 cm.
                \item SAM 3 gera máscara (seleção central).
                \item Área mínima 0,3 m²; confiança 0,25.
                \item Polígonos integrados ao mapa final.
            \end{itemize}
        \end{column}
    \end{columns}
\end{frame}

\begin{frame}{Métricas e Resultados (Abordagem Neural)}
    \begin{columns}[T]
        \begin{column}{0.48\textwidth}
            \textbf{Contagem final}
            \begin{itemize}
                \item 144 parcelas; 60 com CP plantada.
                \item 2.256 CP plantadas.
                \item 1.226 CP mapeadas.
                \item Taxa global: 54,3\%.
                \item 1.228 copas segmentadas.
            \end{itemize}
        \end{column}
        \begin{column}{0.48\textwidth}
            \textbf{Categorias}
            \begin{itemize}
                \item YOLO + SAM 3: 564 (46,0\%).
                \item Só YOLO: 267 (21,8\%).
                \item Só SAM 3: 395 (32,2\%).
                \item Área média: 5,55 m².
                \item Área mediana: 4,29 m².
            \end{itemize}
        \end{column}
    \end{columns}

    \vspace{0.4em}
    \begin{block}{Leitura por composição}
        Melhores taxas: D12 (68,3\%), D7 (65,9\%), D8 (65,7\%) e D11 (62,7\%).
        A D1 concentrou o maior desafio: 1.200 plantadas e 441 mapeadas (36,8\%).
    \end{block}

    \footnotesize Pendente: métricas de teste do detector (P/R/F1, mAP50 e
    mAP50-95), curvas e figuras qualitativas.
\end{frame}

% ═════════════════════════════════════════════════════════════════════════════
\partdivider{Síntese}{Comparação e Conclusões}
% ═════════════════════════════════════════════════════════════════════════════

\begin{frame}{Comparação das Duas Abordagens}
    \begin{table}
        \centering
        \small
        \begin{tabular}{lll}
            \toprule
            \textbf{Aspecto} & \textbf{Clássica (HSV)} & \textbf{Neural (SAM 3 + YOLOv8)} \\
            \midrule
            Sinal & Cor + forma & Padrões aprendidos \\
            Anotação & Não requer & Curadoria por exclusão \\
            Interpretabilidade & Alta & Baixa \\
            Custo & Baixo ($\sim$400 ms/tile, CPU) & Alto (GPU, 500 épocas) \\
            Segmentação & Fecho convexo & Máscara SAM 3 \\
            Discriminação & Fraca (espectral) & Potencialmente alta \\
            Resultado & P 86,5 / R 69,4 / F1 77,0 & 1.226 mapeadas; 54,3\% global \\
            \bottomrule
        \end{tabular}
    \end{table}
    \vspace{0.3em}
    \small As frentes se complementam: a clássica como \textit{baseline} interpretável; a neural buscando precisão e generalização.
\end{frame}

\begin{frame}{Conclusões e Trabalhos Futuros}
    \begin{itemize}
        \item O pipeline \textbf{clássico} é rápido, interpretável e adequado a uma
        primeira estimativa espacial: 2.003 regiões, 86,5\% de precisão, 69,4\% de recall.
        \item O pipeline \textbf{neural} reduz a anotação via SAM 3 e evita vazamento com
        divisão espacial; mapeou 1.226 indivíduos e segmentou 1.228 copas.
        \item \textbf{Próximos passos (clássica)}: refinar HSV, consolidar detecções entre
        tiles, testar descritores de textura.
        \item \textbf{Próximos passos (neural)}: incorporar P/R/F1 e mAP, revisar
        composições sem CP plantada e testar replicabilidade para outras espécies.
        \item \textbf{Integração}: clássica como pré-filtro rápido, neural como refinamento.
    \end{itemize}
\end{frame}

\end{document}
